% =====================================================================================
%  W33 READER BOXES -- the shared audience system, Pass 4316.
%
%  The machine blueprint had a three-audience convention (cream = plain language,
%  blue = specification, rose = withdrawn claim) and the other manuscripts had none:
%  476 pages of w33_paper and 346 of photonic_holonet with no way in for a reader who
%  is not already inside the project.
%
%  NAMING, learned the expensive way.  The first version of this file defined
%  environments called `plain`, `spec` and `gotwrong` to match the blueprint.  That
%  broke photonic_holonet catastrophically -- 106 errors, 10 pages -- because that
%  manuscript uses \spec as a MATH OPERATOR (\operatorname{spec}) declared by its
%  inserts via \providecommand.  Defining a `spec` ENVIRONMENT in the preamble defines
%  the macro \spec, so the later \providecommand silently declined and every \spec(...)
%  in the body invoked an environment opener instead of a spectrum.
%
%  The environments here are therefore named plainbox / specbox / warnbox: distinct from
%  anything the manuscripts already use, and distinct from the blueprint's own names so
%  the blueprint is untouched.
%
%  Colours are the blueprint's, unchanged:
%    plainbg  F4F1EA  cream   -- plain language, for the non-specialist
%    specbg   EAF0F4  blue    -- the exact statement, for the specialist
%    warnbg   F7EDE9  rose    -- a claim this project withdrew
% =====================================================================================

\usepackage{xcolor}
\usepackage{tikz}
\usetikzlibrary{positioning}          % \readerguide places nodes with `right=of`

\providecolor{plainbg}{HTML}{F4F1EA}
\providecolor{specbg}{HTML}{EAF0F4}
\providecolor{warnbg}{HTML}{F7EDE9}

\newenvironment{plainbox}[1][In plain language]
  {\par\medskip\noindent\begin{tikzpicture}
     \node[fill=plainbg,rounded corners=2pt,inner sep=8pt,
           text width=\dimexpr\linewidth-20pt]
     (B) \bgroup\textbf{\small #1.}\itshape\small}
  {\egroup;\end{tikzpicture}\par\medskip}

\newenvironment{specbox}[1][Exactly]
  {\par\medskip\noindent\begin{tikzpicture}
     \node[fill=specbg,rounded corners=2pt,inner sep=8pt,
           text width=\dimexpr\linewidth-20pt]
     (B) \bgroup\textbf{\small #1.}\small}
  {\egroup;\end{tikzpicture}\par\medskip}

\newenvironment{warnbox}[1][What we got wrong]
  {\par\medskip\noindent\begin{tikzpicture}
     \node[fill=warnbg,rounded corners=2pt,inner sep=8pt,
           text width=\dimexpr\linewidth-20pt]
     (B) \bgroup\textbf{\small #1.}\small}
  {\egroup;\end{tikzpicture}\par\medskip}

% The three-audience guide, as one command so every manuscript opens the same way.
\providecommand{\readerguide}[1]{%
  \begin{center}
  \resizebox{\linewidth}{!}{%
  \begin{tikzpicture}[node distance=4mm]
    \node[fill=plainbg,rounded corners,inner sep=7pt,text width=0.32\linewidth,
          font=\small] (rga)
      {\textbf{If you are not a specialist.}\\[3pt]
       Read the cream boxes and the diagrams. They are self-contained: each explains
       the idea before the symbols, and you can skip every blue box without losing
       the argument.};
    \node[fill=specbg,rounded corners,inner sep=7pt,text width=0.32\linewidth,
          font=\small,right=5mm of rga] (rgb)
      {\textbf{If you are a physicist or mathematician.}\\[3pt]
       The blue boxes carry the exact statements, with the scope each one actually
       establishes. Where a result is someone else's, it says so; where it is
       exhaustion rather than proof, it says that too.};
    \node[fill=warnbg,rounded corners,inner sep=7pt,text width=0.32\linewidth,
          font=\small,right=5mm of rgb] (rgc)
      {\textbf{If you are assessing this.}\\[3pt]
       The rose boxes are claims this project published and then withdrew, each with
       the measurement that overturned it. They are the reason the remaining numbers
       can be trusted. #1};
  \end{tikzpicture}}
  \end{center}}
